\documentclass[oneside, landscape]{article}

\usepackage{listings}
\usepackage{color}
\usepackage[letterpaper, scale=0.89, centering]{geometry}
\usepackage{fancyhdr}
\usepackage{hyperref}
\setlength{\parindent}{0em}
\setlength{\parskip}{1em}
\usepackage{menukeys}
\usepackage{emoji}

\usepackage{fontspec}
\directlua{luaotfload.add_fallback
   ("emojifallback",
    {
      "NotoColorEmoji:mode=harf;"
    }
   )}

% https://stackoverflow.com/questions/4439605/c-source-code-in-latex-document
\lstset{
  language=C,                     % choose the language of the code
  numbers=left,                   % where to put the line-numbers
  stepnumber=1,                   % the step between two line-numbers.
  numbersep=5pt,                  % how far the line-numbers are from the code
  backgroundcolor=\color{white},  % choose the background color. You must add \usepackage{color}
  showspaces=false,               % show spaces adding particular underscores
  showstringspaces=false,         % underline spaces within strings
  showtabs=false,                 % show tabs within strings adding particular underscores
  tabsize=2,                      % sets default tabsize to 2 spaces
  breaklines=true,                % sets automatic line breaking
  breakatwhitespace=true,         % sets if automatic breaks should only happen at whitespace
  basicstyle=\ttfamily\footnotesize,
  upquote=true,
  escapeinside={(*@}{@*)},
}

\lstnewenvironment{shell}[1][]
  {\lstset{frame=single, linewidth=8cm, numbers=none, #1}}
  {}

\pagestyle{fancy}
\fancyhf{}
\renewcommand{\headrulewidth}{0pt}

\makeatletter
\lst@InputCatcodes
\def\lst@DefEC{%
 \lst@CCECUse \lst@ProcessLetter
  ^^80^^81^^82^^83^^84^^85^^86^^87^^88^^89^^8a^^8b^^8c^^8d^^8e^^8f%
  ^^90^^91^^92^^93^^94^^95^^96^^97^^98^^99^^9a^^9b^^9c^^9d^^9e^^9f%
  ^^a0^^a1^^a2^^a3^^a4^^a5^^a6^^a7^^a8^^a9^^aa^^ab^^ac^^ad^^ae^^af%
  ^^b0^^b1^^b2^^b3^^b4^^b5^^b6^^b7^^b8^^b9^^ba^^bb^^bc^^bd^^be^^bf%
  ^^c0^^c1^^c2^^c3^^c4^^c5^^c6^^c7^^c8^^c9^^ca^^cb^^cc^^cd^^ce^^cf%
  ^^d0^^d1^^d2^^d3^^d4^^d5^^d6^^d7^^d8^^d9^^da^^db^^dc^^dd^^de^^df%
  ^^e0^^e1^^e2^^e3^^e4^^e5^^e6^^e7^^e8^^e9^^ea^^eb^^ec^^ed^^ee^^ef%
  ^^f0^^f1^^f2^^f3^^f4^^f5^^f6^^f7^^f8^^f9^^fa^^fb^^fc^^fd^^fe^^ff%
  ^^^^c3a9%
  ^^^^fffd%
  ^^00}
\lst@RestoreCatcodes


\lstdefinestyle{shell}{frame=single, linewidth=34em, numbers=none, basicstyle=\ttfamily\footnotesize}

\begin{document}

\small

%% PAGE 1: Review

\begin{minipage}[t]{0.45\textwidth}

\textbf{Inside \texttt{start\_server}}

\begin{lstlisting}[frame=single, numbers=none, basicstyle=\ttfamily\footnotesize, language={}, linewidth=35em]
// http-server.c
typedef void (*RequestHandler)(char *, int);
void start_server(RequestHandler handler, int port) {

  char buffer[2048];

  ... set up on the given port ...

  while(1) {
      // not the real system call names
      int client = accept_connection_from_internet();
      read_from_internet(buffer, client);
      buffer[bytes] = '\0';
      (*handler)(buffer, client);
      close(client);
  }
}
\end{lstlisting}

\begin{lstlisting}[frame=single, basicstyle=\ttfamily\footnotesize, numbers=none, linewidth=35em]
// string-server.c
void handle(char *request, int client) { /* all of PA5 */  }

int main() { start_server(&handle, 2900); }
\end{lstlisting}

\end{minipage}
\hfill
\begin{minipage}[t]{0.50\textwidth}

\vspace{20em}

\end{minipage}

\vspace{1em}
\hrule
\vspace{1em}

\begin{minipage}[t]{0.45\textwidth}

\textbf{Reference: \texttt{sscanf}}

\vspace{0.5em}

Reads formatted data from a string — the inverse of \texttt{sprintf}.

\begin{lstlisting}[numbers=none, basicstyle=\ttfamily\footnotesize]
char name[50]; int age;
sscanf("Alice 30", "%s %d", name, &age);
// name is "Alice", age is 30
// %s stops at whitespace
\end{lstlisting}

\vspace{0.5em}

\textbf{Reference: \texttt{strstr}}

\vspace{0.5em}

\texttt{strstr(haystack, needle)} returns a pointer to the first occurrence of \texttt{needle} in \texttt{haystack}, or \texttt{NULL}.

\begin{lstlisting}[numbers=none, basicstyle=\ttfamily\footnotesize]
char *p = strstr("key=value", "=");
// p points to "=value"
// p + 1 points to "value"
\end{lstlisting}

\end{minipage}
\hfill
\begin{minipage}[t]{0.50\textwidth}

\textbf{HTTP Request and Response}

\vspace{0.5em}

A request from \texttt{curl "localhost:2900/add?s=hello"}:

\begin{lstlisting}[numbers=none, basicstyle=\ttfamily\footnotesize]
GET /add?s=hello HTTP/1.1\r\n
Host: localhost:2900\r\n
User-Agent: curl/8.7.1\r\n
\end{lstlisting}

\vspace{0.5em}

A response sent back with \texttt{write(client, ...)}:

\begin{lstlisting}[numbers=none, basicstyle=\ttfamily\footnotesize]
HTTP/1.1 200 OK\r\n
Content-Type: text/plain\r\n
\r\n
...response body...
\end{lstlisting}

\end{minipage}

\newpage

%% PAGE 2: Program 4

\begin{minipage}[t]{0.40\textwidth}

\textbf{Program 4: Keeping state across requests}

Goal is to have this kind of behavior:

\lstinputlisting[style=shell, linewidth=20em]{../03-03-server/client-output4.txt}

\vspace{0.5em}

A server runs forever. Data that persists between requests can't live on the
stack — it needs to be in global variables or on the heap.

\vspace{1em}

What if \texttt{add\_string} just stores the pointer?

\begin{lstlisting}[basicstyle=\ttfamily\footnotesize, numbers=none]
void add_string(char *s) {
    strings[num_strings] = s;
    num_strings++;
}
\end{lstlisting}



\end{minipage}
\hfill
\begin{minipage}[t]{0.55\textwidth}

% Lines 4-18: globals + add_string signature and guard
\lstinputlisting[basicstyle=\ttfamily\footnotesize, firstline=4, lastline=18, firstnumber=4, aboveskip=0pt, belowskip=0pt]{../03-03-server/string-server4.c}%
\vspace{4\baselineskip}% blank space for lines 19-20 (malloc + strcpy)
% Lines 21-26: rest of add_string + respond_with_list header + for loop
\lstinputlisting[basicstyle=\ttfamily\footnotesize, firstline=21, lastline=26, firstnumber=21, aboveskip=0pt, belowskip=0pt]{../03-03-server/string-server4.c}%
\vspace{2\baselineskip}% blank space for line 27 (write call)
% Lines 28-50: rest of file
\lstinputlisting[basicstyle=\ttfamily\footnotesize, firstline=28, firstnumber=28, aboveskip=0pt, belowskip=0pt]{../03-03-server/string-server4.c}

\end{minipage}

\newpage

%% PAGE 3: Review Questions

\textbf{Review Questions}

\vspace{1em}

1. Given the following request string:

\begin{lstlisting}[numbers=none, basicstyle=\ttfamily\footnotesize]
char request[] = "GET /add?s=goodbye HTTP/1.1\r\nHost: localhost:2900\r\n";
\end{lstlisting}

Write code that extracts \texttt{"goodbye"} into a variable called
\texttt{value}. Use any functions you want, but also refer to {\tt sscanf} and
{\tt strstr} from the handout if needed!

\vspace{3em}

2. Given the following request string:

\begin{lstlisting}[numbers=none, basicstyle=\ttfamily\footnotesize]
char request[] = "GET /add-number?n=367 HTTP/1.1\r\nHost: localhost:2900\r\n";
\end{lstlisting}

Write code that extracts the number \texttt{367} into an \texttt{int} variable called \texttt{n}.

\vspace{3em}

3. Our http server gives an \lstinline{int client} value to the handler. What
is that value used for?



\newpage

%% PAGE 4: Function pointer comparison (instructor reference)

\begin{minipage}[t]{0.45\textwidth}

\textbf{Function Pointers: Java vs C}

\vspace{0.5em}

In Java, you pass behavior using an \textit{interface}:

\begin{lstlisting}[frame=single, language=Java, basicstyle=\ttfamily\footnotesize, numbers=none]
// http-server.java
interface RequestHandler {
  void handle(String request, int client);
}

class Server {
  static void startServer(RequestHandler handler, int port) {
    // server implementation
    ... handler.handle(requestFromClient, clientFileDescriptor); ...
  }
}
\end{lstlisting}

\begin{lstlisting}[frame=single, language=Java, basicstyle=\ttfamily\footnotesize, numbers=none]
// string-server.java
class StringHandler implements RequestHandler {
    public void handle(String request, int client) { /* ... */  }
}
class Main {
  public static void main(String[] args) {
    Server.startServer(new StringHandler(), 2900);
  }
}
\end{lstlisting}

\end{minipage}
\hfill
\begin{minipage}[t]{0.50\textwidth}

\vspace{2em}

In C, there are no interfaces. We pass a \textit{function pointer} — the
address of a function. We can use a \lstinline{typedef} to name the type to
make it look more like Java:

\begin{lstlisting}[frame=single, basicstyle=\ttfamily\footnotesize, numbers=none]
// http-server.c
typedef void (*RequestHandler)(char *, int);

void start_server(RequestHandler handler, int port) {
    // server implementation
    ... (*handler)(request_from_client, client_file_descriptor); ...
}
\end{lstlisting}

\begin{lstlisting}[frame=single, basicstyle=\ttfamily\footnotesize, numbers=none]
// string-server.c
void handle(char *request, int client) { /* ... */  }

int main() { start_server(&handle, 2900); }
\end{lstlisting}

\vspace{0.5em}

In Python, functions can be passed directly, which is more like C:

\begin{lstlisting}[frame=single, language=Python, basicstyle=\ttfamily\footnotesize, numbers=none]
# http-server.py
def start_server(handler, port):
    ... handler(request_from_client, client_file_descriptor) ...

# string-server.py
def handle(request, client): ...

start_server(handle, 2900)
\end{lstlisting}

\end{minipage}

\end{document}
